\documentclass[11pt]{article}
\usepackage[utf8]{inputenc}
\usepackage{amsmath,amssymb}
\usepackage[margin=1in]{geometry}
\usepackage{hyperref}
\emergencystretch=3em

\title{The one-dimensional Taylor tree to all orders, monomial perturbation}
\author{Claude, with Daniel Murfet}
\date{2026-09-06}

\begin{document}
\maketitle
\sloppy

\begin{abstract}
The general-order version of units~28--29: for $\xi(u)=\xi_0+cu^m$ and every truncation order $N$,
\[
\begin{aligned}
&\Bigl|Z_{\rm chart}-\sum_{p<N}\tfrac{(\beta c)^p}{p!}\tfrac1{2k}n^{-\mu_p}S_{\mu_p+p/2}(\xi_0)\Bigr|\\
&\qquad\le\tfrac{(\beta|c|)^N}{N!}\tfrac1{2k}n^{-\mu_N}S_{\mu_N+N/2}(\xi_0+|c|b^m)
+\sum_{p<N}\tfrac{(\beta|c|)^p}{p!}\,\text{tail}_p,\qquad \mu_p=\tfrac{h+mp+1}{2k},
\end{aligned}
\]
with the tails exponentially small. This is \texttt{thm:TaylorTree}/\texttt{cor:standardintegralexp}
to every order in one dimension: the exponents $\mu_p$ run along the arithmetic progression
$(h{+}1)/(2k)+(m/2k)\mathbb N\subseteq\Lambda(h,k)$ and every coefficient is an explicit fluctuation
function. The analytic input is a global remainder bound $|e^x-\sum_{p<N}x^p/p!|\le\frac{|x|^N}{N!}e^{|x|}$.
Zero \texttt{sorry}/\texttt{axiom}.
\end{abstract}

\section{Setting}
Unit~28 handled $N=2$ by hand. The same decomposition $e^x=\sum_{p<N}x^p/p!+R_N(x)$ with
$x=\beta\sqrt nu^kcu^m$ works at every order: the finite Taylor part is a sum of insertion integrals
(each known exactly up to a tail by units~22/24/25), and the remainder is dominated pointwise by
$\frac{|x|^N}{N!}e^{|x|}$, which on the chart is $\frac{(\beta|c|)^N}{N!}$ times the $N$-th insertion
integrand at the shifted argument $\xi_0+|c|b^m$ --- so the remainder integral is again an insertion
value, at order $\mu_N+N/2$.

\section{The things formalised}
{\small
\begin{verbatim}
theorem abs_exp_sub_sum_range_le (x : ℝ) (N : ℕ) :
    |Real.exp x - ∑ p ∈ Finset.range N, x ^ p / p.factorial|
      ≤ |x| ^ N / N.factorial * Real.exp |x|

theorem standardIntegral1D_taylor_tree (β n ξ₀ c b : ℝ) (h k m N : ℕ)
    (hβ : 0 < β) (hn : 1 ≤ n) (hb : 0 < b) (hk : 0 < k) :
    |Z_chart - ∑ p ∈ range N, (β c)^p/p! * ((1/(2k)) n^(-μ_p) S_{μ_p+p/2}(ξ₀))|
      ≤ (β|c|)^N/N! * ((1/(2k)) n^(-μ_N) S_{μ_N+N/2}(ξ₀ + |c| b^m))
        + ∑ p ∈ range N, (β|c|)^p/p! * tail_p
\end{verbatim}
}

\section{Proof strategy}
The exponential bound is the unit-28 series argument with the shift $2$ replaced by $N$; the only new
ingredient is $N!\,p!\le(p+N)!$, from \texttt{Nat.factorial\_mul\_factorial\_dvd\_factorial\_add}. The
main theorem decomposes the integrand pointwise as $\sum_{p<N}G_p+g_R$ with
$G_p(u)=\frac{(\beta c)^p}{p!}u^{h+mp}(\sqrt nu^k)^p e^{\rm pop}(u)$ (the identity
$G_p(u)=u^he^{\rm pop}(u)\,x^p/p!$ plus \texttt{Finset.mul\_sum}), integrates term by term
(\texttt{integral\_finsetSum}, \texttt{integral\_add}; all pieces continuous on the chart), rewrites each
$\int G_p$ as insertion value minus tail exactly as in unit~26, bounds $|\int g_R|$ through the
dominating $N$-th insertion at $\xi_0+|c|b^m$, and closes with $|\sum|\le\sum|\cdot|$ on the tails.

\section{Roadblocks and resolutions}
Essentially none --- the file compiled at the first attempt apart from a deprecated name
(\texttt{continuous\_finset\_sum} $\to$ \texttt{continuous\_finsetSum}). The design choices that made this
smooth were all inherited from units~26 and~28: opaque local definitions via
\texttt{obtain ⟨g, hg⟩ : ∃ g, g = \dots}, a universally quantified algebraic
\texttt{key : ∀ SA ST IR, SA - ST + IR - SA = IR - ST} for the final rearrangement (robust to the goal's
exact parenthesisation), and stating the remainder through the insertion lemma rather than by hand
(no \texttt{rpow} merge needed).

\section{What was Mathlib, what was new}
Mathlib: the exponential series, \texttt{hasSum\_nat\_add\_iff'}, \texttt{hasSum\_le},
\texttt{Nat.factorial\_mul\_factorial\_dvd\_factorial\_add}, \texttt{integral\_finsetSum},
\texttt{continuous\_finsetSum}, \texttt{Finset.abs\_sum\_le\_sum\_abs}. New: the global exponential
Taylor remainder bound (a lean-common candidate alongside its $N=2$ case) and the all-orders
one-dimensional Taylor tree with explicit coefficients and remainder.

\section{Lessons}
Once the $N=2$ case is written with opaque local definitions and the insertion identity as the only
evaluator, the general-$N$ proof is a near-verbatim generalisation; do the small case first only if it
teaches the idioms, then generalise immediately rather than stacking special cases.

\section{Follow-ups}
The $O(n^{-\mu_N})$ packaging (absorbing the $N$ tails, each $n^{p/2}e^{-\varepsilon n}$) is the same
routine as unit~29; the multivariate tree and analytic (non-monomial) $\xi$ remain.
\end{document}
